% TOP_PROBLEM: {"id": 3142, "title": "Chromatic number of $\\frac{3}{3}$-power of graph", "category_id": 3, "category": {"id": 3, "name": "graph_theory", "display_name": "Graph Theory", "description": "Problems involving graphs, networks, and their properties.", "slug": "graph-theory", "order_index": 3, "created_at": "2026-07-31T15:26:25.671Z"}, "status": "open", "problem_number": "OPG-60055", "set_id": 12}
% FIRST_POSTED: 2026-10-01
% ENTRY_KIND: statement_only
% UnsolvedMath ID 3142; source catalogue code OPG-60055
% !TeX program = lualatex
% Definition and sources only; this file is not an LLM solution attempt.
\documentclass[11pt,a4paper]{article}
\usepackage[margin=25mm]{geometry}
\usepackage{fontspec}
\setmainfont{lmroman10-regular.otf}[BoldFont=lmroman10-bold.otf,
 ItalicFont=lmroman10-italic.otf,BoldItalicFont=lmroman10-bolditalic.otf]
\usepackage{amsmath,amssymb,mathtools}
\usepackage{unicode-math}
\setmathfont{latinmodern-math.otf}
\AtBeginDocument{\let\setminus\smallsetminus}
\usepackage{xurl,hyperref}
\hypersetup{colorlinks=true,urlcolor=blue}
\setcounter{secnumdepth}{0}
\setlength{\parindent}{0pt}
\setlength{\parskip}{5pt}
\setlength{\emergencystretch}{3em}
\begin{document}
\section*{Chromatic number of $\frac{3}{3}$-power of graph}\label{3142}
{\small UnsolvedMath ID 3142 · OPG-60055 · Graph Theory · OpenGarden}
\subsection{Definitions and mathematical statement}
All graphs here are finite, simple, and undirected. For an integer
$n\geq1$, the $n$-subdivision $G^{1/n}$ is obtained by replacing
each edge $uv$ of $G$ by a path of exactly $n$ edges from $u$ to
$v$. The $n-1$ internal vertices on each such path are new, and
internal vertices belonging to different original edges are
distinct. For an integer $m\geq1$, the $m$th power of a graph $H$
joins two distinct vertices when their distance in $H$ is at most
$m$. Vertices in different components have infinite distance.

Define the fractional-power notation by
\[
 G^{m/n}=(G^{1/n})^m.
\]
In particular, $G^{3/3}$ means the third power of the graph obtained
by subdividing every edge twice. The fraction in this notation is
not reduced: $G^{3/3}$ need not be $G$.

Write $\Delta(G)$ for the maximum vertex degree of the original
graph. A proper $q$-colouring assigns one of $q$ colours to each
vertex, assigning different colours at adjacent vertices; $\chi$
is the least possible $q$. The conjecture is
\[
 \chi(G^{3/3})\leq2\Delta(G)+1
 \qquad\text{for every }G\text{ with }\Delta(G)\geq2.
\]
Consequently the desired colours must separate every pair within
three steps in the subdivided graph, including both original and
new vertices. The degree in the bound is measured before either
operation.
\subsection{Short English statement}
After replacing each edge by a three-edge path, can all vertices within distance three be distinguished using at most twice the original maximum degree plus one colours?
\subsection{Sources}
\begin{itemize}
\item[S1] ULAM.AI and the UnsolvedMath Contributors, supplied problems.json, record 3142 (OPG-60055), OpenGarden collection. Catalogue identity and source transcription; CC BY 4.0.
\url{https://huggingface.co/datasets/ulamai/UnsolvedMath}.
\item[S2] Open Problem Garden, Chromatic number of $\frac{3}{3}$-power of graph. Original problem and discussion; the mathematical formulation below is expanded from the supplied source record.
\url{http://www.openproblemgarden.org/op/chromatic_number_of_frac_3_3_power_of_graph}.
\end{itemize}
\end{document}
